\documentclass{article}
\usepackage{verbatim}
\usepackage{hyperref}
\usepackage{amssymb}
\usepackage[margin=1.2in]{geometry}

\title{On the Reconciliation of NumPy Headers\\
  in the Construction of GDAL Python Bindings\\[6pt]
  \large A Literate Account of a Build-Order Defect\\
  in OpenDroneMap}

\author{With apologies to D.~E.~Knuth}
\date{April 2026}

\begin{document}
\maketitle

\section{The Problem, Stated Plainly}

A user of OpenDroneMap version 3.6.0, upon processing a set of aerial
photographs, encounters the following fatal message during the report
generation stage:

\begin{verbatim}
  ImportError:
  A module that was compiled using NumPy 1.x cannot be run in
  NumPy 2.3.2 as it may crash.
\end{verbatim}

\noindent The error arises deep within the GDAL library, at the moment when
the hillshade merge script attempts to read raster data into a NumPy
array:

\begin{verbatim}
  File "opendm/tiles/hsv_merge.py", line 198
    rScanline = rBand.ReadAsArray(0, i, ...)
  File "osgeo/gdal.py", line 8910, in ReadAsArray
    from osgeo import gdal_array
  File "osgeo/gdal_array.py", line 10
    from . import _gdal_array
  ImportError: numpy.core.multiarray failed to import
\end{verbatim}

\noindent The processing pipeline then crashes entirely, presenting to the
bewildered user a misleading error about ``strange values in the
reconstruction.''  The photographs are fine; the reconstruction is
fine; the trouble lies elsewhere.

\section{The Architecture of the Build}

To understand the defect, we must first understand how ODM assembles
its dependencies.  The build proceeds in three phases, orchestrated by
\texttt{configure.sh}:

\medskip
\noindent\textbf{Phase 1: \texttt{installreqs}.} System packages are
installed via \texttt{apt-get}, reading the dependency lists from
\texttt{snapcraft24.yaml}.  Among the GDAL build dependencies we find:

\begin{verbatim}
  gdal:
    build-packages:
      - python3-numpy    # <-- Ubuntu 24.04 ships NumPy 1.26.4
      - python3-dev
      - swig
      ...
\end{verbatim}

\noindent A Python virtual environment is then created, and NumPy 2.3.2
is installed into it:

\begin{verbatim}
  python3 -m venv venv --system-site-packages
  venv/bin/pip install numpy==2.3.2
\end{verbatim}

\noindent We now have \emph{two} NumPy installations on the system: version
1.26.4 at \texttt{/usr/lib/python3/dist-packages/numpy/}, installed by
the Debian package manager, and version 2.3.2 at
\texttt{/code/venv/lib/python3.12/site-packages/numpy/}, installed by
pip.

\medskip
\noindent\textbf{Phase 2: SuperBuild.}  CMake compiles GDAL from source
with \texttt{-DBUILD\_PYTHON\_BINDINGS=ON}.  GDAL's build system uses
SWIG to generate Python wrappers, then invokes \texttt{setuptools} to
compile the C extension module \texttt{\_gdal\_array.so}.  The
\texttt{setuptools} build calls \texttt{numpy.get\_include()} to
locate the NumPy C headers.

Here is the crux of the matter: \emph{which} NumPy does
\texttt{setuptools} find?

Despite our CMake arguments pointing to the virtual environment's
Python, GDAL's internal build machinery discovers the \emph{system}
Python's NumPy~1.26.4.  The C extension \texttt{\_gdal\_array.so} is
thus compiled against NumPy~1.x ABI headers and installed to:

\begin{verbatim}
  /code/SuperBuild/install/local/lib/python3.12/
    dist-packages/osgeo/_gdal_array.cpython-312-...-linux-gnu.so
\end{verbatim}

\medskip
\noindent\textbf{Phase 3: \texttt{installpython}.}  The pip
requirements are installed into the virtual environment, including
\texttt{numpy==2.3.2}.  At runtime, the \texttt{PYTHONPATH} places the
SuperBuild directory (containing the NumPy~1.x-compiled
\texttt{\_gdal\_array.so}) \emph{before} the virtual environment's
\texttt{site-packages}.  Thus when any code calls
\texttt{ReadAsArray()}, the import of \texttt{\_gdal\_array} loads a
binary incompatible with the running NumPy, and the process crashes.

\section{Why the Na\"{\i}ve CMake Fix Was Insufficient}

Our first attempt was to pass explicit Python and NumPy paths to
GDAL's CMake configuration, following the pattern already established
by OpenCV and PDALPython in their respective \texttt{External-*.cmake}
files:

\begin{verbatim}
  if (WIN32)
    set(GDAL_PYTHON_ARGS
      -DPython3_EXECUTABLE=${PYTHON_EXE_PATH}
      -DPython3_ROOT_DIR=${PYTHON_HOME}
      -DPython3_NumPy_INCLUDE_DIRS=...numpy/_core/include)
  else()
    set(GDAL_PYTHON_ARGS
      -DPython3_EXECUTABLE=${PYTHON_EXE_PATH}
      -DPython3_NumPy_INCLUDE_DIRS=...numpy/_core/include)
  endif()
\end{verbatim}

\noindent This is good practice and we retain it in our final patch, but it
proved insufficient.  The CMake \texttt{FindPython3} module honors
\texttt{Python3\_EXECUTABLE} during the \emph{configure} step, but GDAL's
SWIG binding compilation occurs later via \texttt{setuptools}, which
performs its own NumPy discovery by importing numpy in whatever Python
it finds.  The system NumPy~1.x, sitting at a well-known path, wins
the race.

\section{The Effective Fix}

The solution is disarmingly simple: before the SuperBuild runs, ensure
that the system NumPy is upgraded to the same major version we intend
to use at runtime.  In \texttt{configure.sh}, immediately after the
virtual environment's NumPy is installed:

\begin{verbatim}
  # edt requires numpy to build
  venv/bin/pip install numpy==2.3.2

  # Upgrade system numpy to 2.x so GDAL Python bindings
  # (_gdal_array.so) are compiled against NumPy 2.x headers.
  sudo pip3 install --break-system-packages \
                    --ignore-installed numpy==2.3.2
\end{verbatim}

\noindent Two flags deserve explanation:

\begin{itemize}
  \item \texttt{-{}-break-system-packages} permits pip to install into
    the system Python's \texttt{dist-packages}, which is normally
    protected on modern Debian-derived systems (PEP~668).

  \item \texttt{-{}-ignore-installed} is necessary because the
    Debian-packaged NumPy~1.26.4 has no pip \texttt{RECORD} file; without
    this flag, pip attempts to uninstall the existing version, fails to
    find the manifest, and aborts.
\end{itemize}

\noindent After this line executes, both the system Python and the virtual
environment agree: NumPy is version~2.3.2.  When GDAL's
\texttt{setuptools} build calls \texttt{numpy.get\_include()}, it
receives the NumPy~2.x header path regardless of which Python
interpreter it uses.  The resulting \texttt{\_gdal\_array.so} is
ABI-compatible with the runtime NumPy, and the import succeeds.

\section{A Defensive Measure}

Even with the root cause eliminated, we observed a secondary defect in
the error cascade.  When the hillshade merge script failed, the
function \texttt{generate\_colored\_hillshade()} returned \texttt{None}
for the output path.  The calling code in \texttt{odm\_report.py}
passed this \texttt{None} directly to \texttt{gdal\_translate} as a
filename:

\begin{verbatim}
  gdal_translate -outsize 1400 0 -of png "None" ...
  ERROR 4: None: No such file or directory
\end{verbatim}

\noindent This converted a recoverable warning (no hillshade preview in
the PDF report) into a fatal pipeline crash.  We add a guard:

\begin{verbatim}
  colored_dem, hillshade_dem, colored_hillshade_dem = \
      generate_colored_hillshade(resized_dem_file)

  if colored_hillshade_dem is not None \
          and os.path.isfile(colored_hillshade_dem):
      system.run("gdal_translate ...")
  else:
      log.ODM_WARNING(
        "Colored hillshade not generated, "
        "skipping DEM preview")

  for f in [resized_dem_file, colored_dem,
            hillshade_dem, colored_hillshade_dem]:
      if f is not None and os.path.isfile(f):
          os.remove(f)
\end{verbatim}

\noindent Now, should any future defect prevent hillshade generation, the
pipeline continues gracefully.  The PDF report will lack a DEM
preview image, but the user's hours of photogrammetric processing
will not be lost.

\section{Summary of Changes}

Three files were modified:

\begin{enumerate}
  \item \textbf{\texttt{configure.sh}} --- Upgrade system NumPy to 2.3.2
    before the SuperBuild, so GDAL's \texttt{\_gdal\_array.so} compiles
    against NumPy~2.x headers.  \emph{This is the essential fix.}

  \item \textbf{\texttt{SuperBuild/cmake/External-GDAL.cmake}} --- Pass
    \texttt{Python3\_EXECUTABLE} and
    \texttt{Python3\_NumPy\_INCLUDE\_DIRS} to GDAL's CMake, matching the
    convention used by OpenCV and PDALPython.  \emph{Belt to the
    suspenders.}

  \item \textbf{\texttt{stages/odm\_report.py}} --- Check for
    \texttt{None} before invoking \texttt{gdal\_translate} on the
    hillshade output.  \emph{A guard against future misfortune.}
\end{enumerate}

\section{Moral}

When a build system has two installations of the same library---one
placed by the system package manager, one by the language package
manager---the build will invariably find the wrong one.  The only
reliable remedy is to ensure they agree.

\bigskip
\hfill\textit{``Beware of bugs in the above code;}

\hfill\textit{I have only proved it correct, not tried it.''}

\hfill--- D.~E.~Knuth (1977)

\end{document}
